\section{Rational descent and four residue-disk orbits}
\label{sec-fixed-rational-descent}

This section concerns the single smooth projective curve
\[
 C:\quad W^2=z^6-27z^4+99z^2-9
\]
and the two minimal elliptic models and actual quotient maps
\[
\begin{aligned}
 E_1 &: y^2=x^3-9x-9, &R_1&=((z^2-9)/4,W/8),\\
 E_2 &: y^2=x^3-189x+999, &R_2&=((33-9/z^2)/4,-9W/(8z^3)).
\end{aligned}
\]
The notation $F_0,T_i,\delta_i,\Xi,\alpha_i^0,R_{0i},L_i$ is that
of the preceding zero-slope expansion. Put
$f=F_0+(4/3)\log_5 3$. All statements here are ordinary proofs,
with the established removable extensions on $C(\mathbb Q_5)$.

\paragraph{RD1: an actual multiplication formula.}
On either short model put $\psi=\psi_9$, $\phi=x\psi_9^2-\psi_{10}\psi_8$,
and $D=\phi'\psi-2\phi\psi'$. These are integral polynomials in $x$:
the two even-index factors contribute $(2y)^2=4(x^3+ax+b)$.
Then the rational-function identities are
\[
 x([9]Q)=\frac\phi{\psi^2},\qquad
 y([9]Q)=\frac{yD}{9\psi^3},\qquad
 T(Q)=-\frac{9\phi\psi}{yD},\qquad
 \delta(Q)=-\frac{9\phi}{yD}=-\frac\phi{\omega_9}.
\]
Indeed the first is the division-polynomial multiplication formula
\cite{SutherlandDivisionPolynomials2023}. If $X=\phi/\psi^2$, the identity
$[9]^*(dx/(2y))=9\,dx/(2y)$ gives
$X'(x)/(2y([9]Q))=9/(2y)$ and $X'=D/\psi^3$.
This fixes both the factor nine and the ordinate sign. The remaining
identities follow from $t=-x/y$. They extend as rational identities;
$D\ne0$ because $X$ is nonconstant in characteristic zero.

\paragraph{RD2: descent to $u=z^2$.}
Write $S(u)=u^3-27u^2+99u-9$, $x_1=(u-9)/4$, $x_2=(33-9/u)/4$.
Define rational functions over $\mathbb Q$ by
\[
\begin{split}
 \mathcal X(u)&=-9u^{39}\frac{\phi_1(x_1)D_2(x_2)}{\phi_2(x_2)D_1(x_1)},\\
 V_1(u)&=\frac{5184\phi_1(x_1)^2\psi_1(x_1)^2}{S(u)D_1(x_1)^2},\\
 V_2(u)&=\frac{64u^3\phi_2(x_2)^2\psi_2(x_2)^2}{S(u)D_2(x_2)^2}.
\end{split}
\]
For actual local points $\Xi=\mathcal X(z^2)$ and $T_i^2=V_i(z^2)$.
To check the constants, $y_2/y_1=-9/z^3$ makes the delta ratio times
$z^{81}$ equal to the first expression, while
$y_1^2=S/64$, $y_2^2=81S/(64u^3)$ give the other two.
The identities on the common open chart extend to the previously
proved removable values. Thus $\mathcal X$ is finite and nonzero and
$V_i\in25\mathbb Z_5$ on the image of the actual local curve.
These assertions are not made for arbitrary $u\in\mathbb Q_5$.

\paragraph{RD3: a square-root-free expression.}
There are unique rational series $\widehat R_i,\widehat M_i$ with
$\widehat R_i(t^2)=R_{0i}(t)$ and $\widehat M_i(t^2)=L_i(t)^2$.
On the actual curve image,
\[
\begin{split}
 F_0={}&-\frac2{81}\left(\log_5\mathcal X(u)
          +\widehat R_1(V_1(u))-\widehat R_2(V_2(u))\right)\\
 &-\frac{\alpha_1^0}{81}\widehat M_1(V_1(u))
   +\frac{\alpha_2^0}{81}\widehat M_2(V_2(u)).
\end{split}
\]
This follows by regrouping the convergent even series and the convergent
Cauchy product $L_i^2$. The proved tail bounds apply at the original
$t$-degrees. For finite nonzero residue classes, $2z$ is a unit and
$u$ is a local coordinate. At zero it has ramification order two;
at infinity $1/u=q^2$ has order two. Algebraic cancellation alone does
not supply analytic denominator or logarithm error bounds.

\paragraph{DS1: exact symmetry.}
The commuting involutions $\iota(z,W)=(z,-W)$ and $\sigma(z,W)=(-z,W)$
generate a group of order four preserving $f$. Indeed
\[
 R_1\iota=-R_1,\quad R_2\iota=-R_2,\quad
 R_1\sigma=R_1,\quad R_2\sigma=-R_2.
\]
Since $\psi_9$ depends only on $x$, $\delta(-Q)=-\delta(Q)$.
The two delta signs cancel under $\iota$; under $\sigma$ the signs
of $z^{81}$ and $\delta_2$ cancel. Thus $\Xi$ is invariant.
The remaining terms are unchanged because $R_{0i}$ is even and
$L_i^2$ is even. Analytic continuation proves invariance at the
exceptional fibers too.

\paragraph{DS2: complete coverage.}
There are twelve smooth residue disks and exactly four orbits:
two over $z=0$, four over $z=\pm1$, four over $z=\pm2$, and two at
infinity. Modulo five the sextic takes value one at zero and four
at every nonzero residue, with nonzero ordinates. The monic leading
coefficient supplies the two infinities. The involutions give precisely
the asserted orbits. Representatives can be chosen through
$(0,\sqrt{-9})$, $(1,8)$, $(2,\sqrt{19})$, and $W/z^3=1$ at infinity.
Their finite parameter $z-a$ or infinity parameter $q=1/z$ ranges over
$5\mathbb Z_5$. Transport sends the parameter to itself or its negative,
so is an analytic isomorphism and preserves each zero multiplicity.
A complete certificate on each representative therefore covers all disks.

\paragraph{DS3: even orders at the fixed fibers.}
At zero the unique Hensel branch satisfies $W(-z)=W(z)$, and $\sigma$
preserves the representative disk. At infinity let $v=Wq^3$; its equation
is $v^2=1-27q^2+99q^4-9q^6$. The composite $\sigma\iota$ preserves
$v(0)=1$ and sends $q$ to $-q$. Hence $f$ is exactly even in both
parameters. Its central derivative vanishes. If its central value is
zero and the series is nonzero, its order there is positive and even.
The following disk certificates determine those values and orders.
